\subsection{The Gamma Distribution}

\begin{definition}
    The Gamma function is defined for $t > 0$ as
    \[\Gamma(t) = \int_0^{\infty} x^{t - 1} e^{-x} dx.\]


\end{definition}
\begin{lemma}
    \[\Gamma(t + 1) = \int_0^{\infty} x^t e{-x} dx \]
    \[ = t \Gamma(t).\]
    \[= t (t - 1) \Gamma(t - 1). \quad \text{ and so on }\]
\end{lemma}
\begin{proof}
    Use integration by parts.
\end{proof}

We can calculate that $\Gamma(1) = 1$, and hence we can show that for positive integers
\[\Gamma(t) = (t - 1) !.\]

\begin{definition}
    The random variable $X$ has a Gamma distribution with parameters $\lambda$ and $k$ $(\lambda > 0, k > 0)$, denoted $Gamma(\lambda, k)$, if it has PDF
    \[f_X(x) = \left\{
        \begin{array}{lr}
          \frac{1}{\Gamma(k)} \lambda^kx^{k - 1} e^{-\lambda x} & x \geq 0\\
          0 & \text{otherwise}
        \end{array}
    \right. .\]

    The Gamma Distribution has $\mathbb{E}[X] = \frac{k}{\lambda}$ and $Var[X] = \frac{k}{\lambda^2}$

    The $\Gamma(\lambda, 1)$ distribution is equal to the $Exp(\lambda)$ distribution.

\end{definition}

\begin{prop}
    If $X \sim \Gamma(\lambda, k)$ and $Y = cX$, then
    \[Y \sim \Gamma(\frac{\lambda}{c}, k).\]
\end{prop}

\begin{prop}
    If $X_1,...,X_k$ are independent variables where $X_i \sim Exp(\lambda) \forall i$, then
    \[\sum_{i=1}^k X_i \sim \Gamma(\lambda, k).\]

    Hence for integers $k_1$ and $k_2$, if $Y_1 \sim \Gamma(\lambda, k_1)$ and $Y_2 \sim \Gamma(\lambda, k_2)$ then
    \[Y_1 + Y_2 \sim \Gamma(\lambda, k_1 + k_2).\]
\end{prop}

\begin{prop} $ $

  \begin{itemize}
    \item If $X \sim N(0,1)$, then $X^2 \sim \Gamma( \frac{1}{2}, \frac{1}{2})$.
    \item If $X_1,...,X_k$ are independent $N(0,1)$ random variables, then
      \[\sum_{i=1}^k X_i^2 \sim \Gamma( \frac{1}{2}, \frac{k}{2}).\]
  \end{itemize}

  The second distribution is also known as the $\chi^2$ distribution on $k$ degrees of freedom.
    
\end{prop}
